\section{Conclusiones}

Un Blueprint de Arquitectura de Seguridad reúne los principales elementos que deben considerarse para
responder a las necesidades de seguridad de una organización, precisa el propósito de cada uno y
muestra cómo se articulan. El caso del Centro de Control Integrado (CCI) de TransMilenio~S.A. permite
recorrer esos elementos en una aplicación práctica, desde la visión, los principios y el alcance de la
arquitectura hasta el análisis de la situación actual, la identificación de brechas, el diseño de la
arquitectura objetivo y la planificación de su implementación.

El análisis del CCI mostró que la seguridad es crítica en un entorno que centraliza procesos de
supervisión y regulación de la flota. Se identificaron riesgos relacionados con la manipulación de la
telemetría, los accesos no autorizados, la trazabilidad de las decisiones operativas y la
disponibilidad de los servicios. Aunque existen controles como la autenticación, las herramientas de
seguridad perimetral y el almacenamiento inmutable WORM, su integración, cobertura y efectividad aún
pueden mejorar.

La arquitectura objetivo (TO-BE) se organiza en cuatro dominios principales: gobierno de identidades
y control de acceso, integridad de la telemetría, auditoría y trazabilidad de las decisiones, y
continuidad operativa. Los dominios organizan los controles según los riesgos identificados y
resumen las capacidades que necesita el CCI para operar de forma confiable, controlada y resiliente.

Las arquitecturas de transición y la hoja de ruta llevan esa visión a una implementación progresiva,
organizada en horizontes de tiempo. Esa secuencia permite priorizar las acciones: primero el
fortalecimiento de los mecanismos de acceso, después la validación e integridad de la información, y
al final las capacidades de alta disponibilidad y recuperación ante incidentes.

SABSA y O-ESA, junto con los modelos y las vistas de ArchiMate, permiten relacionar las necesidades
del negocio con los riesgos, los controles, los servicios y los componentes tecnológicos de la
arquitectura. Los indicadores y las métricas, por su parte, sirven para evaluar el cumplimiento de
los objetivos de seguridad y hacer seguimiento a la efectividad de las iniciativas propuestas.

El caso del CCI muestra que un Blueprint de Arquitectura de Seguridad es una guía que conecta las
necesidades de la organización con los riesgos, las capacidades de seguridad y las acciones
necesarias para alcanzar un estado objetivo. Los controles y las tecnologías se eligen dentro de esa
guía, después de definir los riesgos y las capacidades. Su elaboración exige considerar en conjunto
los aspectos estratégicos, organizacionales, tecnológicos y operativos, y traducir el diseño en
decisiones y procesos de implementación.
